% Experiment 2: POSIX File Copy

\section{Experiment 2: POSIX File Copy with open/read/write}

\subsection{Problem Statement}
The objective is to implement a CLI tool named \texttt{fdcopy} that copies a byte stream using low-level file descriptor I/O.

The tool must:
\begin{itemize}
    \item Accept inputs: \verb|--src <path>| and \verb|--dst <path>| (required).
    \item Optional inputs: \verb|--buf <N>| and \verb|--force|.
    \item Read from STDIN if source is \verb|-|.
    \item Compute the IEEE 802.3 CRC32 checksum of copied bytes.
    \item Prevent overwriting destination unless \verb|--force| is used.
    \item Output format: \verb|OK: COPIED <bytes> BYTES| and \verb|OK: CRC32 <hex>|.
\end{itemize}

\subsection{Source Code}
\begin{lstlisting}[language=C, caption={fdcopy.c}]
// ###### PSOIX O/R/W ##########
#include <stdio.h>
#include <stdlib.h>
#include <string.h>
#include <fcntl.h>
#include <unistd.h>
#include <errno.h>
#include <stdint.h>

// Error codes
#define E_USAGE "E_USAGE"
#define E_OPEN_SRC "E_OPEN_SRC"
#define E_OPEN_DST "E_OPEN_DST"
#define E_EXISTS "E_EXISTS"
#define E_READ "E_READ"
#define E_WRITE "E_WRITE"
#define E_CLOSE "E_CLOSE"
#define E_RANGE "E_RANGE"

// Printing Errors
void print_error(const char *code, const char *msg) {
    fprintf(stderr, "ERROR: %s: %s\n", code, msg);
    exit(1);
}

// CRC32 implementation, code outsourced
uint32_t crc32_update(uint32_t crc, const unsigned char *p, size_t len) {
    crc = ~crc;
    while (len--) {
        crc ^= *p++;
        for (int i = 0; i < 8; i++)
            crc = (crc >> 1) ^ ((crc & 1) ? 0xEDB88320 : 0);
    }
    return ~crc;
}

int main(int argc, char *argv[]) {
    char *src_path = NULL;
    char *dst_path = NULL;
    int buffer_size = 4096;
    int force = 0;

    // Parsing arguments
    for (int i = 1; i < argc; i++) {
        if (strcmp(argv[i], "--src") == 0) {
            if (i + 1 < argc) src_path = argv[++i];
        } else if (strcmp(argv[i], "--dst") == 0) {
            if (i + 1 < argc) dst_path = argv[++i];
        } else if (strcmp(argv[i], "--buf") == 0) {
            if (i + 1 < argc) {
                buffer_size = atoi(argv[++i]);
                if (buffer_size < 1 || buffer_size > 1048576) {
                    print_error(E_RANGE, "buffer size 1..1048576");
                }
            }
        } else if (strcmp(argv[i], "--force") == 0) {
            force = 1;
        }
    }

    if (!src_path || !dst_path) {
        print_error(E_USAGE, "missing --src or --dst");
    }

// Opening src file
    int fd_in;
    if (strcmp(src_path, "-") == 0) {
        fd_in = STDIN_FILENO;
    } else {
        fd_in = open(src_path, O_RDONLY);
        if (fd_in < 0) print_error(E_OPEN_SRC, strerror(errno));
    }

    // Opening dest file
    int flags = O_WRONLY | O_CREAT;
    if (!force) {
        flags |= O_EXCL;
    } else {
        flags |= O_TRUNC;
    }

    int fd_out = open(dst_path, flags, 0644);
    if (fd_out < 0) {
        if (errno == EEXIST) {
            print_error(E_EXISTS, "destination already exists (use --force)");
        } else {
            print_error(E_OPEN_DST, strerror(errno));
        }
    }

    // Copy loop
    unsigned char *buf = malloc(buffer_size);
    if (!buf) print_error(E_USAGE, "memory allocation failed");

    ssize_t n_read;
    size_t total_bytes = 0;
    uint32_t crc = 0;

    // Read until end of file(EOF)
    while ((n_read = read(fd_in, buf, buffer_size)) > 0) {
        crc = crc32_update(crc, buf, n_read);
        
        char *ptr = (char*)buf;
        ssize_t n_write_total = 0;
        
        // Partial writes
        while (n_write_total < n_read) {
            ssize_t res = write(fd_out, ptr + n_write_total, n_read - n_write_total);
            if (res < 0) print_error(E_WRITE, strerror(errno));
            n_write_total += res;
        }
        total_bytes += n_read;
    }

    if (n_read < 0) print_error(E_READ, strerror(errno));

    // Cleaning
    free(buf);
    if (fd_in != STDIN_FILENO) close(fd_in);
    if (close(fd_out) < 0) print_error(E_CLOSE, strerror(errno));

    // Output
    printf("OK: COPIED %zu BYTES\n", total_bytes);
    printf("OK: CRC32 %08x\n", crc);

    return 0;
}
\end{lstlisting}

\subsection{Sample Input and Output}

\textbf{Case 1:}
\begin{verbatim}
$ ./fdcopy --src source.txt --dst dest.txt
OK: COPIED 3 BYTES
OK: CRC32 352441c2
\end{verbatim}

\textbf{Case 2}
\begin{verbatim}
$ ./fdcopy --src - --dst out/q02.bin
ERROR: E_EXISTS: destination already exists (use --force)
\end{verbatim}